% recon_otb_v2 8pt-minimum version (2026-11-02). Source: simspect repo paper_figures/drafts/recon_otb_v2_8pt.tex; 7pt version: recon_otb_v2.tex.
\begin{figure}[t]
\centering
% 8pt-minimum variant of recon_otb_v2 (all text \footnotesize, register/memory digits full size: $r0$ not $r0$).%
% Results figure: the Recon-STT getOldestTaint (OTB) bug, hand-built test, real gem5 pipeview.%
% Body only (goes into the paper's tikz/). Wrapped in { }; helper names are prefixed \ro / ro.%
% Needs TikZ libraries calc + shapes.geometric (the paper preamble loads both).%
% Listing = oldest_taint_bug/debug_out/inst-oldest-taint-bug.xml (xml_readable.txt): idx IX0..IX5 =%
%   i0..i5; Reg_s$0 = r0, Reg_s$1 = r1, Mem_s$0 = m0; kinds, flags, rf and ddi edges as in the XML.%
%   Omitted: the XML's rf i1 -> i2 (the hand-built asm replaces both branch conditions with slow%
%   pointer-chase loads, run.py custom_pass25; inst-oldest-taint-bug.s) and i5's inmem m0.%
%   i4's input order (r1, r0) is a drawing choice; the XML's two inputs are unordered.%
% Trace = oldest_taint_bug/debug_out/pipeview_scheme2.txt (/work/gem5-recon --scheme=2,%
%   wiki/results/reconresults.md "Bug 1"). Cycle c = (tick - 2803000)/500:%
%   i0 = seq 562 (jne uop .2) complete 2803000 -> 0; i5 = seq 575 issue 2807000 -> 8 (complete 0,%
%   retire 0: no completion recorded before the squash); i2 = seq 571 complete 2849000 -> 92.%
% Mechanism: taint_tracker.cc getOldestTaint 119-131, propagateTaints 108-116, freeTaints 134-155%
%   (reconresults.md "Source references"). Plan + caveats: paper_figures/plans/RESULTS_FIGURES.md.%
{%
\providecolor{ssspecfill}{gray}{0.94}\providecolor{ssspecink}{gray}{0}%
\providecolor{ssprotfill}{gray}{0.83}\providecolor{ssprotink}{gray}{0}%
\providecolor{ssxmitfill}{gray}{0.76}\providecolor{ssxmitink}{gray}{0}%
\providecolor{sssrcfill}{gray}{0.66}\providecolor{sssrcink}{gray}{0}\providecolor{sssrcbg}{gray}{0.90}%
% ============================ TUNING KNOBS ============================%
\def\roshade{sssrcbg}%  defense source of truth (Recon's taint rule box: light bg + sssrcink outline)%
\def\rokind{13pt}%       x where kinds start (ids at 0)%
\def\roout{61pt}%        x of the output operand%
\def\roarr{72pt}%        x of the <- arrow%
\def\roin{84pt}%         x of the first input operand%
\def\rogut{116pt}%       x of the gutter carrying i1's rf edge%
\def\roflag{119pt}%      x where flag words start%
\def\ronote{162pt}%      x where notes start%
\def\ronotew{88pt}%      note width (\ronote + \ronotew <= 252pt)%
\def\rorow{15pt}%        default row pitch%
\def\rogA{9pt}%          extra space above i1 (its ddi arc)%
\def\rogC{9pt}%          extra space above i3 (its ddi arc)%
\def\rogD{14pt}%         extra space above i4 (incoming rf over its ddi arc)%
\def\rogE{14pt}%         extra space above i5 (i4's lower ddi arc)%
\def\rolw{1.1pt}%        line width of the leak path%
\def\rocyc{11}%          pt per cycle on the timeline%
\def\rotx{64}%           pt, x of cycle -2 (timeline start)%
\def\rotxb{212}%         pt, x of cycle 90 (after the elision)%
\def\rorowh{17}%         pt between timeline rows%
\def\robarh{4}%          pt height of a bar%
\def\rostar{8.5pt}%      outer diameter of the transmission star%
\def\rowinlw{0.9pt}%     line width of the window frame%
% ======================================================================%
\newcommand\rox[1]{(\rotx+(#1+2)*\rocyc)}% cycles -2..10%
\newcommand\roxb[1]{(\rotxb+(#1-90)*\rocyc)}% cycles 90..93%
\begin{tikzpicture}[x=1pt, y=1pt, >=latex,%
  every node/.style={font=\footnotesize, inner sep=0pt},%
  execute at begin node={\hyphenpenalty=10000\exhyphenpenalty=10000},%
  rotok/.style={anchor=base west, outer sep=0.6pt},%
  rolab/.style={font=\footnotesize\itshape},%
  roflagw/.style={font=\footnotesize\itshape, anchor=base west},%
  ronote/.style={font=\footnotesize\itshape, anchor=base west, align=left, text width=\ronotew},%
  roleak/.style={->, line width=\rolw, draw=ssprotink, preaction={draw=ssprotfill, arrows=-, line width=3.6pt, line cap=round, blend mode=multiply}},%
  rorowl/.style={font=\footnotesize, anchor=base west},%
]%
\path (0,0) coordinate (y0) (0,-\rorow-\rogA) coordinate (y1) (0,-2*\rorow-\rogA) coordinate (y2)%
  (0,-3*\rorow-\rogA-\rogC) coordinate (y3) (0,-4*\rorow-\rogA-\rogC-\rogD) coordinate (y4)%
  (0,-5*\rorow-\rogA-\rogC-\rogD-\rogE) coordinate (y5);%
\path (0,8pt);% breathing room above%
% speculative rows (after the unresolved i2): a light band behind ids/kinds/operands only, ending before the flag column so no colored text sits on it; drawn first%
\fill[ssspecfill, rounded corners=2pt] ($(-2pt,0 |- y2)+(0,-4.5pt)$) rectangle ($(\ronote-4pt,0 |- y5)+(0,-4pt)$);%
\coordinate (spA) at ($(-2pt,0 |- y2)+(0,-4.5pt)$); \coordinate (spB) at ($(\ronote-4pt,0 |- y5)+(0,-4pt)$);%
\node[font=\footnotesize\itshape, text=black!55, rotate=90, anchor=south] at ($($(spB)!0.5!(spB |- spA)$)+(-1.2pt,0)$) {speculative};% vertical label along the tan band's right edge
% ------------------------------ listing ------------------------------%
\foreach \r in {0,...,5} \node[rotok] at (y\r) {i\r:};%
% i0: Branch_x  committed%
\node[rotok] at ($(y0)+(\rokind,0)$) {\texttt{Branch\_x}};%
\node[roflagw] at ($(y0)+(\roflag,0)$) {committed};%
% i1: Load r0 <- m0  committed  (load A, older)%
\node[rotok] at ($(y1)+(\rokind,0)$) {\texttt{Load}};%
\node[rotok] (i1o) at ($(y1)+(\roout,0)$) {$r0$};%
\node[rotok] at ($(y1)+(\roarr,0)$) {$\leftarrow$};%
\node[rotok] (i1m) at ($(y1)+(\roin,0)$) {$m0$};%
\node[roflagw] at ($(y1)+(\roflag,0)$) {committed};%
% i2: Branch_x  (unresolved)%
\node[rotok] at ($(y2)+(\rokind,0)$) {\texttt{Branch\_x}};%
% i3: Load r1 <- m0  (load B, younger)%
\node[rotok] at ($(y3)+(\rokind,0)$) {\texttt{Load}};%
\node[rotok] (i3o) at ($(y3)+(\roout,0)$) {$r1$};%
\node[rotok] at ($(y3)+(\roarr,0)$) {$\leftarrow$};%
\node[rotok] (i3m) at ($(y3)+(\roin,0)$) {$m0$};%
% i4: Other_n r0 <- r1, r0  (the merge)%
\node[rotok] at ($(y4)+(\rokind,0)$) {\texttt{Other\_n}};%
\node[rotok] (i4o) at ($(y4)+(\roout,0)$) {$r0$};%
\node[rotok] at ($(y4)+(\roarr,0)$) {$\leftarrow$};%
\node[rotok] (i4a) at ($(y4)+(\roin,0)$) {$r1$};%
\node[rotok] (i4c) at (i4a.base east) {,};%
\node[rotok] (i4b) at ($(i4c.base east)+(2.5pt,0)$) {$r0$};%
% i5: Load r0 <- [r0]  xmit%
\node[rotok] at ($(y5)+(\rokind,0)$) {\texttt{Load}};%
\node[rotok] at ($(y5)+(\roout,0)$) {$r0$};%
\node[rotok] at ($(y5)+(\roarr,0)$) {$\leftarrow$};%
\node[rotok] (i5b) at ($(y5)+(\roin,0)$) {[};%
\node[rotok, outer sep=0pt] (i5r) at (i5b.base east) {$r0$};%
\node[rotok] at (i5r.base east) {]};%
\node[roflagw, text=ssxmitink] at ($(y5)+(\roflag,0)$) {xmit};%
% ---- edges: the older source (thin) ----%
\draw[->, shorten >=1.2pt] ($(i1m.north)+(0,0.3pt)$) to[out=115, in=65, looseness=1.1]%
  node[pos=0.5, above=1.4pt, rolab] {ddi} ($(i1o.north)+(0,0.6pt)$);%
\draw[->, shorten <=0.6pt, shorten >=0.8pt, rounded corners=2pt]%
  (i1o.south) |- ($(\rogut,0 |- y2)+(0,2.5pt)$) |- (i4b.east);%
\node[rolab, anchor=base west] at ($(\roin,0 |- y2)+(4pt,4.5pt)$) {rf};%
\draw[->, shorten <=0.4pt, shorten >=1pt] ($(i4b.south)+(0,-0.2pt)$) .. controls +(0,-7pt) and +(4pt,-7pt) ..%
  ($(i4o.south east)+(-0.8pt,-0.2pt)$);%
\node[rolab, anchor=north west] at ($(i4b.south)+(1.5pt,-2.5pt)$) {ddi};%
% ---- edges: the younger source = the leak path (thick) ----%
\draw[roleak, shorten >=1.4pt] ($(i3m.north)+(0,0.3pt)$) to[out=115, in=65, looseness=1.1]%
  node[pos=0.5, above=1.4pt, rolab] {ddi} ($(i3o.north)+(0,0.6pt)$);%
\draw[roleak, shorten <=0.6pt, shorten >=1pt] (i3o.south) .. controls +(0.15,-0.22) and +(0.08,0.22) .. ($(i4a.north east)+(-0.6pt,0.4pt)$)%
  node[pos=0.5, rolab, anchor=south west, xshift=2pt] {rf};%
\draw[roleak, shorten >=1.2pt] ($(i4a.north west)+(0.8pt,0.3pt)$) to[out=125, in=55, looseness=0.9]%
  ($(i4o.north)+(0.5pt,0.6pt)$);%
\node[rolab, anchor=base east] at ($(i4o.base west)+(-1pt,5.5pt)$) {ddi};%
\draw[roleak, shorten <=0.6pt, shorten >=1.2pt] ($(i4o.south west)+(0.8pt,0)$) .. controls +(-1pt,-14pt) and +(-8pt,5pt) .. (i5b.north west)%
  node[pos=0.3, rolab, anchor=east, xshift=-4pt] {rf};%
% ---- notes ----%
\node[ronote] at ($(y0)+(\ronote,0)$) {resolves first};%
\node[ronote] at ($(y1)+(\ronote,0)$) {older tainted source};%
\node[ronote] at ($(y2)+(\ronote,0)$) {\textcolor{ssspecink}{mispredicted, resolves last}};%
\node[ronote] at ($(y3)+(\ronote,0)$) {younger tainted source};%
\node[ronote] at ($(y4)+(\ronote,0)$) {keeps only i1's taint};%
% (pipeline timeline removed 2026-11-03, Anna: "fix it or take it out"; trace facts kept in the caption)%
% ------------- the defense-specific rule -------------%
\node[draw=sssrcink, fill=\roshade, anchor=north west, name=rorule, align=left, inner sep=3pt, text width=244pt, font=\footnotesize] at ($(0,0 |- y5)+(0,-14pt)$) {\raggedright\textit{Recon-STT rule:} an ALU result keeps the taint of only its oldest tainted source (\texttt{getOldestTaint}), and the taint is freed once that source is no longer speculative (\texttt{freeTaints}). i4 keeps i1's taint, which is freed when i0 resolves, although i3 still feeds i4 under the unresolved i2.};%
\path ($(rorule.south)+(0,-6pt)$);% breathing room below%
\end{tikzpicture}%
}%
\caption{\claudetag{[TODO: edit caption]} \textbf{Recon-STT releases a transmitter whose address still depends on a speculative load} (\texttt{getOldestTaint}; \texttt{/work/gem5-recon}, STT scheme, hand-built test). Top: the test as an execution graph in the notation of Table~\ref{tab:execvocab} \claudetag{[TODO: Claude fixed this table ref]}. The ALU op i4 merges $r_0$ from the older load i1, speculative under branch i0, with $r_1$ from the younger load i3, speculative under the mispredicted branch i2; its result is the address of transmitter i5 (thick arrows: the flow from i3). In the gem5 trace, i5 issues 8 cycles after i0 resolves, as soon as i4's result is ready and in the same cycle as with no defense, while i2 stays unresolved for 92 cycles.
\claudetag{[TODO: Claude removed the pipeline timeline from this figure and kept its numbers here]} Recon gives i4's result only the taint of its oldest tainted source and frees it when that source stops being speculative (blue), so the dependence on i3 is lost. This test was built by hand; our bounded enumeration does not reach this shape within its instance cap.}
\label{fig:otb}
\end{figure}
